% Part: incompleteness
% Chapter: theories-computability
% Section: computably-axiomatizable

\documentclass[../../../include/open-logic-section]{subfiles}

\begin{document}

\olfileid{inc}{tcp}{cax}

\olsection{可公理化理论}

理论 $\Th{T}$ 被称为是\emph{可公理化的}，如果它有一个可计算的公理集~$A$。（称 $A$ 是 $\Th{T}$ 的公理集，意思是$T = \Setabs{!A}{A \Proves !A}$。）自然数的任何“合理的”公理化都具有这一性质。特别地，任何具有有限公理集的理论都是可公理化的。

\begin{lem}
假设 $\Th{T}$ 是可公理化的。那么 $\Th{T}$ 是可计算枚举的。
\end{lem}

\begin{proof}
假设 $A$ 是 $\Th{T}$ 的一个可计算公理集。要判定 $!A \in T$，只需从公理出发搜索 $!A$ 的一个推导。

换一种说法，$!A$ 属于 $\Th{T}$ 当且仅当存在 $A$ 中公理的一个有限列表 $!B_1$, \dots, $!B_k$，以及一阶逻辑中 $(!B_1 \land \dots \land !B_k) \lif !A$ 的一个推导。但我们已知，任何具有“存在……使得……”形式定义的集合，只要第二个“……”是可计算的，就是 c.e.\@ 的。
\end{proof}

\end{document}